\label{re:review_2}

\begin{center}
  \large\bf Response to Reviewer 2
\end{center}

We would like to thank you for your valuable comments and suggestions. More importantly, we are very pleased that we can have so many beneficial discussions on multisource remote sensing image classification with you. We believe these discussions must be (and it has been) helpful for us to improve this work.

We have very carefully followed all the comments and cleared up every issue raised. We provide below a point-by-point response to the comments as follows.

\begin{tcolorbox}[title={Reviewer \#2, Comment \#1}]
  \vspace{1em}

  \color{blue}
  Has the author considered the issue of class imbalance? From the experimental results, such a problem indeed exists. What are the underlying causes of this issue?

  \vspace{1em}
\end{tcolorbox}

\textbf{Response:}

We sincerely thank the reviewer for this helpful comment. We fully agree that class imbalance is indeed present in our experimental results, and we have carefully analyzed its underlying causes and added a corresponding discussion to the revised manuscript.

Specifically, we would like to clarify the following:

\textbf{(1) Regarding the underlying cause of class imbalance:}

The class imbalance observed in our results primarily stems from the intrinsic, naturally imbalanced land-cover distribution of the benchmark datasets themselves, rather than from a limitation specific to our network design. In both the Augsburg and Houston2018 datasets, classes such as Forest, Water, and Residential area occupy substantially larger and more contiguous spatial extents, and are therefore represented by far more labeled training pixels, whereas classes such as Commercial area and Industrial area correspond to sparse, scattered regions with comparatively few labeled samples. This inherent imbalance in the number of available training samples per class is a common characteristic of real-world land-cover classification benchmarks, and it directly limits the amount of supervisory signal available for learning discriminative representations for minority classes.

\textbf{(2) Regarding a compounding factor beyond sample quantity:}

Beyond the quantitative imbalance in sample counts, we further observe that the classes most affected (e.g., Commercial area) also exhibit strong spectral and structural similarity to neighboring majority classes (e.g., Residential area, Industrial area), since all correspond to man-made built-up surfaces with overlapping reflectance and geometric characteristics. This compounding effect of limited supervision and inherent spectral ambiguity makes these minority classes particularly challenging to classify accurately, as discussed in detail in our Failure Case Analysis (Section III-C) and in our response to Comment~6 of Reviewer 1.

We have added the following clarification on the underlying cause of class imbalance to \textbf{Section III-C} of the revised manuscript, immediately preceding the Failure Case Analysis paragraph already added in response to Comment~6 of Reviewer~1. The specific added text is as follows:

``\textcolor{red}{Although SGFNet is not specifically designed for class-imbalance learning, it partially alleviates the negative effects of imbalanced samples. SGFNet achieves notable improvements on minority classes (Roads, Sidewalks, Crosswalks, and Residential buildings), together with the highest AA of 96.71\%. It indicates that its semantic feature extraction and cross-modal fusion enhance minority-class discrimination. However, without balanced sampling, loss reweighting, or data augmentation, SGFNet mitigates rather than fundamentally solves the class-imbalance problem.}''

We hope this analysis effectively clarifies the underlying causes of class imbalance in our experiments and addresses the reviewer's concern. All corresponding changes have been highlighted in red in the updated manuscript for your review. Thank you again for guiding us to improve our manuscript.

\vspace{2em}

\begin{tcolorbox}[title={Reviewer \#2, Comment \#2}]
  \vspace{1em}
  
  \color{blue}
  How does SGFNet address the issue of spatial misalignment in multi-source data?

  \vspace{1em}
\end{tcolorbox}

\textbf{Response:}

We sincerely thank the reviewer for this important question. We agree that the mechanism by which SGFNet addresses spatial misalignment deserves a more explicit and detailed explanation. To address this concern, we have elaborated on the design rationale of the Frequency Modulated Fusion Block (FMFB) and provided empirical verification through a controlled experiment.

Specifically, we would like to clarify the following:

\textbf{(1) Regarding the mechanism by which FMFB addresses spatial misalignment:}

SGFNet addresses the issue of spatial misalignment primarily through the proposed FMFB, which performs cross-modal feature interaction in the frequency domain rather than directly conducting spatial-domain fusion. The key idea is that, according to the shift theorem of the discrete transform, a spatial translation of an image alters only the phase component of its frequency representation, while the magnitude/energy distribution across frequency bands remains largely unchanged. Since slight cross-modal misalignment in multi-source remote sensing data typically manifests as small spatial translations or local geometric distortions, frequency-domain operations that primarily act on the energy distribution, as performed by FMFB, are inherently less sensitive to such shifts than pixel-wise spatial-domain fusion, which requires precise spatial correspondence between modalities to be effective. Concretely, FMFB first transforms the HSI and SAR/LiDAR features into the frequency domain via 2D DCT, then computes their frequency-domain difference to obtain a joint representation that is used to guide a cross-modal attention mechanism (Eqs.~8--12), and finally transforms the modulated features back into the spatial domain via inverse DCT. This design allows the network to adaptively emphasize reliable, well-aligned frequency components while suppressing the influence of misaligned, high-frequency structural discrepancies, thereby alleviating the adverse impact of spatial misalignment on feature fusion.

\textbf{(2) Regarding empirical verification:}

To empirically verify that FMFB indeed improves robustness to spatial misalignment, we conducted a controlled experiment on the Houston2018 dataset by artificially applying spatial shifts of $\{1, 2, 3, 5\}$ pixels along both the $x$- and $y$-directions to the LiDAR branch input, and compared SGFNet against two representative state-of-the-art methods, DFFNet and MSFMamba, under identical shift conditions. 

The results confirm that SGFNet exhibits substantially smaller performance degradation as the shift magnitude increases compared to DFFNet and MSFMamba, verifying that FMFB effectively addresses the issue of spatial misalignment through frequency-domain modulation. The newly added experiments are as follows:

``\textbf{\textcolor{red}{Robustness to Spatial Misalignment.}} \textcolor{red}{We conduct experiments on the Houston2018 dataset by applying spatial shifts of $\{1, 2, 3, 5\}$ pixels along both the $x$- and $y$-directions to the LiDAR branch input, while keeping the HSI branch fixed. As shown in Table V, the OA of DFFNet and MSFMamba degrades by 2.60\% and 2.03\%, respectively, whereas SGFNet exhibits only 1.40\% drop. It confirms that the proposed FMFB alleviates cross-modal spatial misalignment through frequency-domain modulation.}''

\begin{table}[h]
  \centering
  \footnotesize
  \caption*{\textcolor{red}{\footnotesize TABLE V: Comparison of OA (\%) under different degrees of artificial spatial shifts on the Houston2018 dataset.}}
  \begin{tabular}{lccccc}
  \toprule
  Shift (pixels) & 0 & 1 & 2 & 3 & 5 \\
  \midrule
  DFFNet & 91.21 & 90.98 & 90.09 & 89.36 & 88.61 \\
  MSFMamba & 92.38 & 92.14 & 92.07 & 91.81 & 90.35 \\
  \rowcolor{bg} \textbf{SGFNet (Ours)} & \textbf{94.17} & \textbf{94.03} & \textbf{93.96} & \textbf{93.64} & \textbf{92.77} \\
  \bottomrule
  \end{tabular}
\end{table}




We hope this detailed explanation of the underlying mechanism, together with the empirical verification via the controlled shift experiment, effectively addresses the reviewer's question regarding how SGFNet handles spatial misalignment. All corresponding changes have been highlighted in red in the updated manuscript for your review. Thank you again for guiding us to improve our manuscript.

\vspace{2em}

\begin{tcolorbox}[title={Reviewer \#2, Comment \#3}] 
  \vspace{1em}

  \color{blue}
  The paper mentions that SMCB can better capture long-range semantic dependencies compared to traditional convolution, and how this can be explained through physical meaning.
  
  \vspace{1em}
\end{tcolorbox}

\textbf{Response:}

We sincerely thank the reviewer for this insightful question. We fully agree that the physical meaning underlying SMCB's ability to capture long-range semantic dependencies deserves a more explicit explanation. We have added a detailed clarification to the revised manuscript, which we also briefly summarized in our response to Comment~1 of Reviewer~1.

Specifically, we would like to clarify the following:

\textbf{(1) Regarding why traditional convolution is limited in modeling long-range dependencies:}

Standard convolutional kernels are spatially shared and operate only within a small, fixed local window (e.g., $3\times3$) at every spatial position. Consequently, a single convolutional layer can only aggregate information from its immediate neighborhood, and information from distant, semantically related regions can only be incorporated indirectly by stacking many layers or through pooling operations, which progressively lose fine-grained spatial detail. This makes it difficult for traditional convolution to directly model long-range semantic correlations, particularly in remote sensing scenes where semantically related land-cover regions (e.g., disjoint patches of the same building complex or road network) may be spatially distant from one another.

\textbf{(2) Regarding the physical meaning of how SMCB captures long-range dependencies:}

In contrast, SMCB generates its dynamic convolution kernel by computing an affinity between a per-position query $\mathbf{Q}$ and a set of key vectors $\mathbf{K}$ that are obtained via global adaptive pooling over the entire feature map (Eq.~2), rather than from a local neighborhood. Physically, each of the $K^2$ key vectors is computed by adaptively pooling contextual information aggregated from across the whole spatial extent of the input, and therefore represents a global contextual descriptor rather than a locally restricted one. When computing the affinity map (Eq.~3), every spatial position is allowed to measure its relevance to these globally-derived key vectors, regardless of their physical distance. As a result, the dynamic kernel at a given position can be influenced by, and adaptively aggregate information related to, semantically similar regions located anywhere in the image, not merely its local neighborhood. This is the fundamental physical mechanism through which SMCB achieves long-range semantic dependency modeling: the receptive field of the dynamic kernel is effectively extended from a local spatial window to the semantic structure of the entire feature map, while still being applied through an efficient, spatially localized convolution-like operation (Eq.~6).

We have added the following clarification:

``\textcolor{red}{Physically, the key vectors $\mathbf{K}$ in Eq.~(2) are obtained by adaptively pooling contextual information from across the entire spatial extent of the input feature map, and therefore represent global contextual descriptors rather than locally restricted ones. When computing the affinity map in Eq.~(3), every spatial position measures its relevance to these globally-derived key vectors regardless of their physical distance, allowing the resulting dynamic kernel to adaptively aggregate information from semantically related regions located anywhere in the image. This effectively extends the receptive field of the dynamic kernel from a local spatial window to the semantic structure of the entire feature map, which constitutes the underlying physical mechanism through which SMCB models long-range semantic dependencies.}''

Due to the 5-page limit of IEEE GRSL, we will add these descriptions in the supplementary material. We hope this explanation effectively clarifies the physical meaning underlying SMCB's long-range dependency modeling capability. All corresponding changes have been highlighted in red in the updated manuscript for your review. Thank you again for guiding us to improve our manuscript.

\vspace{2em}

\begin{tcolorbox}[title={Reviewer \#2, Comment \#4}] 
  \vspace{1em}
  
  \color{blue}
  What impact do different convolution kernel sizes in the SMCB module have on the overall detection performance?

  \vspace{1em}
\end{tcolorbox}

\textbf{Response:}

We sincerely thank the reviewer for this valuable suggestion. We fully agree that the sensitivity of SGFNet to the dynamic kernel size $K$ in SMCB is an important parameter that warrants systematic investigation. We have conducted a new sensitivity analysis experiment, which we also refer to in our response to Comment~2 of Reviewer~3, where a similar concern regarding the sensitivity of the dynamic kernel size was raised.

Specifically, we would like to clarify the following:

\textbf{(1) Regarding the experimental setup:}

We evaluated SGFNet with the dynamic kernel size $K$ in SMCB set to $\{5, 7, 9, 11\}$ on both the Augsburg and Houston2018 datasets, while keeping all other hyperparameters (e.g., the number of SGFMs, the number of PCA components) fixed at their default settings described in Section III-B. All other experimental settings follow those described in Section III-A.

\textbf{(2) Regarding the experimental results and analysis:}

As shown in Table~II, the OA on both datasets follows a similar trend: performance improves as $K$ increases from 5 to 7, reaching its peak at $K=7$ (92.58\% on Augsburg and 94.10\% on Houston2018), and then gradually declines as $K$ further increases to 9 and 11. When $K$ is too small (e.g., $K=5$), the number of semantic sampling locations is insufficient to adequately capture the contextual affinities needed for effective semantic modeling, leading to lower OA (91.64\% and 93.26\% on Augsburg and Houston2018, respectively). Conversely, when $K$ is too large (e.g., $K=11$), the dynamic kernel aggregates information from an excessively large and potentially less relevant set of sampling locations, which may introduce noisy or redundant contextual information and increase computational overhead without a corresponding gain in accuracy. Based on this analysis, we set $K=7$ as the default configuration in our experiments, as it achieves the best trade-off between classification accuracy and computational cost.

We have added a new subsection titled \textbf{``Sensitivity of the Dynamic Kernel Size in SMCB''} in \textbf{Section III-B (Parameter Analysis)}, together with a new table presenting the experimental results. The specific added text is as follows:

``\textbf{\textcolor{red}{Sensitivity of the Dynamic Kernel Size in SMCB.}} \textcolor{red}{We evaluate SGFNet with $K \in \{5, 7, 9, 11\}$ on both datasets, and the results are shown in Table \ref{table_kernel}. The OA on both datasets peaks at $K=7$. A smaller $K$ provides insufficient semantic sampling locations for effective contextual modeling, while a larger $K$ introduces potentially redundant or noisy contextual information. Hence, we set $K=7$ in our experiments, as it achieves the best trade-off between classification accuracy and computational overhead.}''

\begin{table}[h]
  \centering
  \footnotesize
  \caption*{\textcolor{red}{\footnotesize TABLE II: OA (\%) of SGFNet under different dynamic kernel sizes $K$ in SMCB.}}
  \begin{tabular}{lcccc}
  \toprule
  $K$ & 5 & 7 & 9 & 11 \\
  \midrule
  Augsburg    & 91.64 & \textbf{92.58} & \underline{92.50} & 92.02 \\
  Houston2018 & 93.26 & \textbf{94.10} & \underline{94.03} & 93.47 \\
  \bottomrule
  \end{tabular}
  \label{table_kernel}
\end{table}

We hope this sensitivity analysis effectively addresses the reviewer's concern regarding the impact of the dynamic kernel size on classification performance. All corresponding changes have been highlighted in red in the updated manuscript for your review. Thank you again for guiding us to improve our manuscript.


\vspace{2em}

\begin{center}
  \rule[-10pt]{14.3cm}{0.05em}
\end{center}

Thanking you again for your precious time and valuable efforts spent in reviewing our manuscript. We greatly appreciate your expertise and the dedicated efforts you have made to assist us in enhancing the quality of this manuscript. We hope that we have clarified all the issues raised by you, and our revised manuscript will again meet your expectations. We eagerly anticipate any additional insights or suggestions you might offer.

Best wishes and thanks again.

\textit{Yuwei Zhao and Feng Gao}

Ocean University of China

Email: gaofeng@ouc.edu.cn
